\section*{Chapitre \ref{ch:g6:latitude-longitude} --- Latitude, longitude et fuseaux horaires}

\begin{solution}{exo:g6:latitude-longitude:1}
«~La \omterm{def:g6:latitude-longitude:latitude}{latitude} se mesure à partir de l'\omterm{def:g3:continents-and-oceans:equator}{équateur} et va de $0^\circ$ à
$90^\circ$ (au nord ou au sud)~; la \omterm{def:g6:latitude-longitude:longitude}{longitude} se mesure à partir du
\omterm{def:g6:latitude-longitude:longitude}{méridien d'origine} et va de $0^\circ$ à $180^\circ$ (à l'est ou à
l'ouest).~»
\end{solution}

\begin{solution}{exo:g6:latitude-longitude:2}
Le \omterm{def:g3:continents-and-oceans:poles}{pôle Nord}~: $90^\circ$\,N. Le \omterm{def:g3:continents-and-oceans:poles}{pôle Sud}~: $90^\circ$\,S. Une ville
située sur l'\omterm{def:g3:continents-and-oceans:equator}{équateur}~: $0^\circ$.
\end{solution}

\begin{solution}{exo:g6:latitude-longitude:3}
Sydney~: \omterm{def:g3:continents-and-oceans:equator}{hémisphère sud}, est. New York~: \omterm{def:g3:continents-and-oceans:equator}{hémisphère nord}, ouest. Le
Caire~: \omterm{def:g3:continents-and-oceans:equator}{hémisphère nord}, est. Rio de Janeiro~: \omterm{def:g3:continents-and-oceans:equator}{hémisphère sud}, ouest.
\end{solution}

\begin{solution}{exo:g6:latitude-longitude:4}
(a) Vrai. (b) Faux~: les \omterm{def:g6:latitude-longitude:latitude}{parallèles} raccourcissent vers les \omterm{def:g3:continents-and-oceans:poles}{pôles}~; le plus
long est l'\omterm{def:g3:continents-and-oceans:equator}{équateur}. (c) Vrai. (d) Faux~: il passe par Greenwich, à
Londres.
\end{solution}

\begin{solution}{exo:g6:latitude-longitude:5}
Le Caire. New York. Tokyo est à environ $36^\circ$\,N, $140^\circ$\,E (un
peu au-dessus du \omterm{def:g6:latitude-longitude:latitude}{parallèle} $30^\circ$\,N, un peu au-delà du \omterm{def:g6:latitude-longitude:longitude}{méridien}
$120^\circ$\,E, avant $150^\circ$\,E).
\end{solution}

\begin{solution}{exo:g6:latitude-longitude:6}
$45 \div 15 = 3$~: leurs midis sont à 3~heures d'intervalle. Il est midi
en premier dans la ville~B, qui est plus à l'est.
\end{solution}

\begin{solution}{exo:g6:latitude-longitude:7}
$73 \div 15 \approx 4.9$, arrondi à 5~: le fuseau est $+5$. Quand il est
midi à Greenwich, il est $12 + 5 = 17$, soit 17~h, à Bombay.
\end{solution}

\begin{solution}{exo:g6:latitude-longitude:8}
$30 \times 111 = 3\,330$~km. Le \omterm{def:g3:continents-and-oceans:poles}{pôle Nord} est à $90^\circ$~:
$90 \times 111 = 9\,990$~km, environ 10\,000~km.
\end{solution}

\begin{solution}{exo:g6:latitude-longitude:9}
$9 - (-5) = 14$~heures d'écart. $8 + 14 = 22$~: il est 22~h le lundi à
Tokyo.
\end{solution}

\begin{solution}{exo:g6:latitude-longitude:10}
Elle a raison. La Terre tourne d'ouest en est, donc les villes de l'est
sont emportées vers le Soleil en premier~: elles le voient se lever,
atteindre son point le plus haut et se coucher avant les villes situées
plus à l'ouest.
\end{solution}

\begin{solution}{exo:g6:latitude-longitude:11}
(a) $42 \times 3 = 126$~secondes, un peu plus de 2~minutes. (b)
$126 \div 240 \approx 0.5$~: environ un demi-degré de \omterm{def:g6:latitude-longitude:longitude}{longitude}. (c)
$0.5 \times 111 \approx 55$~km (avec $0.525$~: environ 58~km). Une erreur
d'une cinquantaine de kilomètres~: assez pour manquer une île, ou pour
s'échouer sur une \omterm{def:g1:landscapes-around-us:water}{côte} que l'on croyait lointaine.
\end{solution}

\begin{solution}{exo:g6:latitude-longitude:12}
(a) $18 - 9 = 9$~: 9~h le lundi à Greenwich. (b) $9 + 10 = 19$~: 19~h le
lundi à Greenwich. (c) $19 - 8 = 11$~: 11~h le lundi à Los Angeles. L'avion
est parti le lundi à 18~h (heure de Tokyo) et atterrit à 11~h le même lundi
(heure de Los Angeles)~: les horloges à l'arrivée indiquent une heure plus
matinale que les horloges au départ, parce que l'avion a traversé la \omterm{def:g6:latitude-longitude:date-line}{ligne
de changement de date} vers l'est et est revenu d'un jour en arrière.
\end{solution}

\begin{solution}{pb:g6:latitude-longitude:1}
\textbf{1.} Ana (Rio de Janeiro, $23^\circ$\,S) et Mia (Auckland,
$37^\circ$\,S).

\textbf{2.} Ana et Sam ($43^\circ$\,O et $118^\circ$\,O). Lena est sur le
\omterm{def:g6:latitude-longitude:longitude}{méridien d'origine}.

\textbf{3.} Lena~: Londres est à $52^\circ$\,N, plus loin de l'\omterm{def:g3:continents-and-oceans:equator}{équateur} que
toutes les autres villes.

\textbf{4.} Mia~: Auckland est à $175^\circ$\,E, à seulement
$180 - 175 = 5$ degrés du \omterm{def:g6:latitude-longitude:longitude}{méridien} $180^\circ$.

\textbf{5.} $360^\circ \div 24 = 15^\circ$.

\textbf{6.} Londres~: $0$. Rio de Janeiro~: $43 \div 15 \approx 2.9$,
ouest, donc $-3$. Tokyo~: $140 \div 15 \approx 9.3$, est, donc $+9$.
Auckland~: $175 \div 15 \approx 11.7$, est, donc $+12$. Los Angeles~:
$118 \div 15 \approx 7.9$, ouest, donc $-8$.

\textbf{7.} Rio de Janeiro~: $20 - 3 = 17$, 17~h le vendredi. Tokyo~:
$20 + 9 = 29$~; $29 - 24 = 5$~: 5~h le samedi.

\textbf{8.} Tokyo est très à l'est~: ses horloges ont 9~heures d'avance
sur celles de Londres. À 20~h à Londres, Tokyo a déjà passé minuit~: son
samedi a donc commencé.

\textbf{9.} $20 - 8 = 12$~: midi le vendredi.

\textbf{10.} $12 - (-8) = 20$~heures d'avance.

\textbf{11.} En traversant la \omterm{def:g6:latitude-longitude:date-line}{ligne de changement de date} vers l'est, le
navire revient au jour précédent.

\textbf{12.} $20 + 12 = 32$~; $32 - 24 = 8$~: quand l'appel commence, il
est \textbf{8~h le samedi} à Auckland. Mia rejoint l'appel à l'heure du
petit déjeuner, pendant que Sam déjeune, le vendredi.
\end{solution}
